\chapter{Matrices}\label{ch:b1:matrices}

A matrix is a \omterm{def:b1:linmaps:def}{linear map} written in \omterm{prop:b1:vspaces:coordinates}{coordinates}. This chapter \omterm{def:b1:logic:sets}{sets}
up the dictionary --- composition becomes matrix product, bijectivity
becomes invertibility, \omterm{def:b1:matrices:changeofbasis}{change of basis} becomes conjugation --- and
the algorithmic side: \omterm{met:b1:matrices:gauss}{row operations}, computation of ranks and
inverses. First met in the High School volume, matrices are now
grounded in the theory of
\cref{ch:b1:vspaces,ch:b1:findim,ch:b1:linmaps}.

\section{Matrices and linear maps}

\begin{definition}\label{def:b1:matrices:def}
$\mathcal{M}_{n,p}(K)$ is the \omterm{def:b1:vspaces:def}{vector space} of $n \times p$ arrays $A
= (a_{ij})$ of scalars ($i$: row, $j$: column), of dimension $np$
(\omterm{def:b1:vspaces:free}{basis}: the matrices $E_{ij}$ with a single $1$). Given bases
$\mathcal{B} = (e_1, \dots, e_p)$ of $E$ and $\mathcal{C}$ of $F$
($\dim F = n$), the \emph{matrix of $u \in \mathcal{L}(E, F)$} is
the array whose $j$-th column lists the \omterm{prop:b1:vspaces:coordinates}{coordinates} of $u(e_j)$ in
$\mathcal{C}$:
\[
\operatorname{Mat}_{\mathcal{B},\mathcal{C}}(u) = (a_{ij}),
\qquad u(e_j) = \sum_{i=1}^{n} a_{ij}\, f_i .
\]
The \omterm{def:b1:logic:map}{map} $u \mapsto \operatorname{Mat}_{\mathcal{B},\mathcal{C}}(u)$
is an isomorphism from $\mathcal{L}(E, F)$ onto
$\mathcal{M}_{n,p}(K)$ (\cref{prop:b1:linmaps:basis}: a \omterm{def:b1:linmaps:def}{linear map}
is exactly a choice of \omterm{def:b1:linmaps:kerim}{images} of the $e_j$).
\end{definition}

\begin{example}[The derivative, as a matrix]
\label{ex:b1:matrices:derivative}
Let $D(P) = P'$ on $\R_3[X]$. In the monomial \omterm{def:b1:vspaces:free}{basis} $(1, X, X^2,
X^3)$: $D(1) = 0$, $D(X) = 1$, $D(X^2) = 2X$, $D(X^3) = 3X^2$, so
\[
\operatorname{Mat}(D) =
\begin{pmatrix}
0 & 1 & 0 & 0\\
0 & 0 & 2 & 0\\
0 & 0 & 0 & 3\\
0 & 0 & 0 & 0
\end{pmatrix}.
\]
In the \emph{divided} \omterm{def:b1:vspaces:free}{basis} $\bigl(1,\ X,\ \frac{X^2}2,\
\frac{X^3}6\bigr)$, each \omterm{def:b1:vspaces:free}{basis} vector \omterm{def:b1:logic:map}{maps} to the previous one
($D\bigl(\frac{X^k}{k!}\bigr) = \frac{X^{k-1}}{(k-1)!}$), and the
matrix becomes the pure shift: ones on the superdiagonal, zeros
elsewhere. Two morals: the matrix belongs to the \emph{pair}
(\omterm{def:b1:logic:map}{map}, \omterm{def:b1:vspaces:free}{basis}), not to the \omterm{def:b1:logic:map}{map} alone; and a good \omterm{def:b1:vspaces:free}{basis} makes
structure visible at a glance --- the shift form shows instantly
that $D^4 = 0$ on $\R_3[X]$, each power of the matrix pushing its
diagonal of ones one step further out.
\end{example}

\begin{definition}[Product]\label{def:b1:matrices:product}
For $A \in \mathcal{M}_{n,p}$ and $B \in \mathcal{M}_{p,q}$:
\[
(AB)_{ik} = \sum_{j=1}^{p} a_{ij}\, b_{jk}
\qquad (1 \leq i \leq n,\ 1 \leq k \leq q).
\]
This is precisely the matrix of the composition:
$\operatorname{Mat}(v \circ u) = \operatorname{Mat}(v)\,
\operatorname{Mat}(u)$ (bases matching in the middle). Similarly,
if $X$ is the column of \omterm{prop:b1:vspaces:coordinates}{coordinates} of $x$, the column of $u(x)$ is
$AX$.
\end{definition}

\begin{proof}[Proof of the composition formula]
\[
v(u(e_k)) = v\Bigl(\sum_j b_{jk} f_j\Bigr) = \sum_j b_{jk}\, v(f_j)
= \sum_j b_{jk} \sum_i a_{ij}\, g_i
= \sum_i \Bigl(\sum_j a_{ij} b_{jk}\Bigr) g_i . \qedhere
\]
\end{proof}

\begin{proposition}[The algebra $\mathcal{M}_n(K)$]\label{prop:b1:matrices:ring}
Square matrices $\mathcal{M}_n(K)$ form a (non-commutative for $n
\geq 2$) \omterm{def:b1:structures:ring}{ring}, with identity $I_n$; its \omterm{def:b1:structures:group}{group} of units is the
\emph{general linear group}\index{general linear group}
$GL_n(K)$, corresponding to \omterm{def:b1:logic:inj}{bijective} endomorphisms. For $A, B \in
\mathcal{M}_n(K)$:
\[
AB = I_n \implies A \in GL_n(K) \text{ and } B = A^{-1}
\]
(one-sided inverses are two-sided, by
\cref{cor:b1:linmaps:samedim}).
\end{proposition}

\begin{proof}
\omterm{def:b1:structures:ring}{Ring} axioms transport from $\mathcal{L}(E)$ through the
isomorphism of \cref{def:b1:matrices:def}: it converts
composition into product and sum into sum, so associativity,
distributivity and the role of $I_n$ are inherited from the
corresponding facts about \omterm{def:b1:logic:map}{maps}, with no entrywise verification.
Non-commutativity: $E_{12}E_{21} = E_{11} \neq E_{22}
= E_{21}E_{12}$. If $AB = I_n$: the endomorphism $a$ of $A$
satisfies $a \circ b = \mathrm{id}$, so $a$ is \omterm{def:b1:logic:inj}{surjective} ($x =
a(b(x))$ exhibits a \omterm{def:b1:logic:map}{preimage} of every $x$), hence \omterm{def:b1:logic:inj}{bijective} in
\omterm{def:b1:findim:def}{finite dimension} (\cref{cor:b1:linmaps:samedim}); composing $a
\circ b = \mathrm{id}$ with $a^{-1}$ on the left gives $b =
a^{-1}$, and then $b\circ a = \mathrm{id}$ too: the one-sided
inverse was two-sided all along --- a strictly
\omterm{def:b1:findim:def}{finite-dimensional} favor.
\end{proof}

\begin{definition}[Transpose; trace]\label{def:b1:matrices:transpose}
The \emph{transpose} of $A = (a_{ij}) \in \mathcal{M}_{n,p}$ is $A^{\mathsf T}
= (a_{ji}) \in \mathcal{M}_{p,n}$; it satisfies $(AB)^{\mathsf T} = B^{\mathsf T}
A^{\mathsf T}$ and $(A^{\mathsf T})^{\mathsf T} = A$. The
\emph{trace}\index{trace} of a square matrix is $\operatorname{tr} A
= \sum_i a_{ii}$; it is \omterm{def:b1:linmaps:def}{linear}, and
\[
\operatorname{tr}(AB) = \operatorname{tr}(BA)
\qquad (A \in \mathcal{M}_{n,p},\ B \in \mathcal{M}_{p,n}).
\]
\end{definition}

\begin{proof}[Proof of the trace identity]
$\operatorname{tr}(AB) = \sum_i \sum_j a_{ij} b_{ji}$ and
$\operatorname{tr}(BA) = \sum_j \sum_i b_{ji} a_{ij}$: the same
double sum.
\end{proof}

\begin{example}[The trace at work]
\label{ex:b1:matrices:tracework}
The \omterm{def:b1:linmaps:projection}{projection} of \cref{ch:b1:linmaps} onto
$\operatorname{Vect}(1,1)$ along $\operatorname{Vect}(0,1)$,
$p(x, y) = (x, x)$, has matrix $A = \begin{pmatrix} 1 & 0\\ 1 &
0\end{pmatrix}$ in the canonical \omterm{def:b1:vspaces:free}{basis}: indeed $A^2 = A$, and
\[
\operatorname{tr} A = 1 = \operatorname{rk} A ,
\]
illustrating \cref{exo:b1:matrices:8}: for idempotents the \omterm{def:b1:matrices:transpose}{trace}
\emph{counts} the dimension of the \omterm{def:b1:linmaps:kerim}{image}, whatever slanted \omterm{def:b1:vspaces:free}{basis}
the matrix is written in. The invariance mechanism is the
identity $\operatorname{tr}(AB) = \operatorname{tr}(BA)$:
\[
\operatorname{tr}\bigl(P^{-1}(AP)\bigr) =
\operatorname{tr}\bigl((AP)P^{-1}\bigr) = \operatorname{tr} A ,
\]
so all matrices similar to $A$ share its \omterm{def:b1:matrices:transpose}{trace} --- the first
\emph{numerical invariant} of an endomorphism, to be joined by
the determinant in \cref{ch:b1:det} (the pair $(s, p)$ of the
weekend problem below).
\end{example}

\begin{example}[Symmetric plus antisymmetric]
\label{ex:b1:matrices:symsplit}
Call $A$ \emph{symmetric} when $A^{\mathsf T} = A$,
\emph{antisymmetric} when $A^{\mathsf T} = -A$. Every square
matrix splits uniquely as one plus the other:
\[
A = \underbrace{\frac{A + A^{\mathsf T}}{2}}_{\text{symmetric}}
+ \underbrace{\frac{A - A^{\mathsf T}}{2}}_{\text{antisymmetric}},
\]
and a matrix that is both is zero ($A = -A$): the two \omterm{def:b1:logic:sets}{sets} are
\omterm{def:b1:vspaces:sum}{supplementary subspaces} of $\mathcal{M}_n(K)$ --- the exact
analogue of the even/odd split of functions
(\cref{ex:b1:vspaces:evenodd}), with transposition playing the
role of $x \mapsto -x$. Dimensions: a symmetric matrix is \omterm{def:b1:vspaces:free}{free}
on and above the diagonal, an antisymmetric one strictly above
(zero diagonal):
\[
\frac{n(n+1)}{2} + \frac{n(n-1)}{2} = n^2 ,
\]
and the count balancing is Grassmann's confirmation of
directness. For $n = 2$: $\begin{pmatrix} 1 & 5\\ 1 &
2\end{pmatrix} = \begin{pmatrix} 1 & 3\\ 3 & 2\end{pmatrix} +
\begin{pmatrix} 0 & 2\\ -2 & 0\end{pmatrix}$. Symmetric
matrices return as the \omterm{def:b1:derivative:def}{second-derivative} data of
\cref{ch:b1:multivar} (the Monge triple $r, s, t$), and the
symmetric-orthogonal ones are classified in
\cref{exo:b1:euclid:12}.
\end{example}

\section{Change of basis}

\begin{definition}\label{def:b1:matrices:changeofbasis}
Let $\mathcal{B}, \mathcal{B}'$ be bases of $E$. The \emph{change of
basis matrix}\index{change of basis} $P =
P_{\mathcal{B}\to\mathcal{B}'}$ has for columns the \omterm{prop:b1:vspaces:coordinates}{coordinates} of
the \emph{new} \omterm{def:b1:vspaces:free}{basis} vectors in the \emph{old} \omterm{def:b1:vspaces:free}{basis}. It is
invertible, $P^{-1} = P_{\mathcal{B}'\to\mathcal{B}}$, and
\omterm{prop:b1:vspaces:coordinates}{coordinates} transform by $X = PX'$ (old $=$ $P\,\cdot$ new).
\end{definition}

\begin{example}[Reading the change-of-basis matrix]
\label{ex:b1:matrices:changerun}
In $\R^2$, from the canonical $\mathcal B$ to $\mathcal B' =
\bigl((1,1), (1,-1)\bigr)$:
\[
P = P_{\mathcal B\to\mathcal B'} =
\begin{pmatrix} 1 & 1\\ 1 & -1 \end{pmatrix}
\]
(new vectors written in old \omterm{prop:b1:vspaces:coordinates}{coordinates}, column by column). The
vector of old \omterm{prop:b1:vspaces:coordinates}{coordinates} $X = (3, 1)^{\mathsf T}$ has new
\omterm{prop:b1:vspaces:coordinates}{coordinates} $X' = P^{-1}X = \frac12(3 + 1,\ 3 - 1)^{\mathsf T} =
(2, 1)^{\mathsf T}$: indeed $2(1,1) + 1(1,-1) = (3,1)$. Mind the
direction --- the matrix $P$ is built from the \emph{new} \omterm{def:b1:vspaces:free}{basis}
but converts \emph{new to old} \omterm{prop:b1:vspaces:coordinates}{coordinates} ($X = PX'$); passing
from old to new costs the inverse. Writing the sanity check
$2(1,1) + (1,-1) = (3,1)$ after every conversion catches the
inverted-$P$ error, which is the most common mistake of the
chapter.
\end{example}

\begin{theorem}[Change of basis for a map]\label{thm:b1:matrices:conjugation}
Let $u \in \mathcal{L}(E)$ with matrix $A$ in $\mathcal{B}$ and $A'$
in $\mathcal{B}'$, and $P = P_{\mathcal{B}\to\mathcal{B}'}$. Then
\[
A' = P^{-1} A\, P .
\]
Two matrices related this way are \emph{similar}\index{similar
matrices}. (For $u \colon E \to F$ with two pairs of bases, the
formula is $A' = Q^{-1} A P$ --- \emph{equivalent}
matrices\index{equivalent matrices}.)
\end{theorem}

\begin{proof}
For any $x$: $X = PX'$ and the \omterm{def:b1:linmaps:kerim}{image} satisfies $Y = AX$, $Y = PY'$.
So $PY' = APX'$, i.e.\ $Y' = (P^{-1}AP)X'$ for all $X'$: the matrix
of $u$ in the new \omterm{def:b1:vspaces:free}{basis} is $P^{-1}AP$ (take for $X'$ the canonical
columns).
\end{proof}

\begin{example}[A good basis makes a map transparent]
\label{ex:b1:matrices:conjugationrun}
Let $u(x, y) = (y, x)$ (swap), with matrix $A = \begin{pmatrix} 0
& 1\\ 1 & 0\end{pmatrix}$ in the canonical \omterm{def:b1:vspaces:free}{basis}. In the \omterm{def:b1:vspaces:free}{basis}
$\mathcal B' = \bigl((1,1), (1,-1)\bigr)$:
\[
P = \begin{pmatrix} 1 & 1\\ 1 & -1 \end{pmatrix},
\qquad
P^{-1} = \frac12\begin{pmatrix} 1 & 1\\ 1 & -1 \end{pmatrix},
\qquad
P^{-1} A P = \begin{pmatrix} 1 & 0\\ 0 & -1 \end{pmatrix}.
\]
No matrix product was really needed: $u$ fixes $(1,1)$ and
reverses $(1,-1)$, so in $\mathcal B'$ its matrix \emph{must} be
$\operatorname{diag}(1, -1)$ --- the swap is the reflection
across the line $y = x$. Finding, for a given endomorphism, a
\omterm{def:b1:vspaces:free}{basis} in which its matrix becomes diagonal is the central problem
of the Year~2 volume (reduction theory); the weekend problem
below shows how far \omterm{def:b1:poly:def}{polynomial} identities alone already go.
\end{example}

\begin{example}[Change of basis, run in reverse]
\label{ex:b1:matrices:reverserun}
The \omterm{def:b1:linmaps:projection}{projection} onto $F = \operatorname{Vect}(1,1)$ along $G =
\operatorname{Vect}(1,-1)$ has, in the adapted \omterm{def:b1:vspaces:free}{basis} $\mathcal
B' = \bigl((1,1),(1,-1)\bigr)$, the transparent matrix $A' =
\operatorname{diag}(1, 0)$. To get its \omterm{def:b1:vspaces:free}{canonical-basis} matrix,
run \cref{thm:b1:matrices:conjugation} backwards, $A = P A'
P^{-1}$:
\[
P = \begin{pmatrix} 1 & 1\\ 1 & -1\end{pmatrix},
\quad
P^{-1} = \frac12\begin{pmatrix} 1 & 1\\ 1 & -1\end{pmatrix},
\quad
A = P\begin{pmatrix} 1 & 0\\ 0 & 0\end{pmatrix}P^{-1}
= \frac12\begin{pmatrix} 1 & 1\\ 1 & 1\end{pmatrix}.
\]
Check: $A^2 = A$ (idempotent), $\operatorname{tr} A = 1 =
\operatorname{rk} A$, and $A\binom{1}{1} = \binom11$,
$A\binom{1}{-1} = 0$, as prescribed. This reverse direction ---
design the matrix in the good \omterm{def:b1:vspaces:free}{basis}, then conjugate back --- is
how rotation, reflection and \omterm{def:b1:linmaps:projection}{projection} matrices are actually
produced in practice.
\end{example}

\begin{theorem}[Rank normal form]\label{thm:b1:matrices:rank}
The \emph{rank} of a matrix (the rank of its columns, equivalently
of the associated \omterm{def:b1:linmaps:def}{linear map}) is the only invariant of equivalence:
every $A \in \mathcal{M}_{n,p}$ of rank $r$ is equivalent to
\[
J_r = \begin{pmatrix} I_r & 0 \\ 0 & 0 \end{pmatrix},
\]
and $\operatorname{rk}(A^{\mathsf T}) = \operatorname{rk}(A)$: row rank
equals column rank.
\end{theorem}

\begin{proof}
Let $u \colon E \to F$ have rank $r$. Choose a \omterm{def:b1:vspaces:sum}{supplementary} $S$ of
$\ker u$ ($\dim S = r$, \cref{thm:b1:linmaps:ranknullity}) with
\omterm{def:b1:vspaces:free}{basis} $(e_1, \dots, e_r)$, completed by a \omterm{def:b1:vspaces:free}{basis} of $\ker u$ into a
\omterm{def:b1:vspaces:free}{basis} of $E$; the \omterm{def:b1:linmaps:kerim}{images} $f_i = u(e_i)$, $i \leq r$, form a \omterm{def:b1:vspaces:free}{basis} of
$\operatorname{im} u$ (the restriction is an isomorphism), completed
into a \omterm{def:b1:vspaces:free}{basis} of $F$. In these bases the matrix of $u$ is exactly
$J_r$. So $A = Q J_r P^{-1}$ for invertible $P, Q$.

Transposing: $A^{\mathsf T} = (P^{-1})^{\mathsf T} J_r^{\mathsf T} Q^{\mathsf T}$
with $J_r^{\mathsf T}$ of the same form (rank $r$) and the outer
factors invertible (\omterm{def:b1:matrices:transpose}{transpose} of invertible is invertible, from
$(AB)^{\mathsf T} = B^{\mathsf T}A^{\mathsf T}$ applied to $AA^{-1} = I$):
$\operatorname{rk} A^{\mathsf T} = r$.
\end{proof}

\section{Row operations}

\begin{method}[Gaussian elimination on matrices]\label{met:b1:matrices:gauss}
The three \emph{elementary row operations}\index{row operations} ---
swap two rows, multiply a row by $\lambda \neq 0$, add a multiple of
a row to another --- do not change the rank (each is left
multiplication by an invertible matrix). Algorithm: create a pivot
(leftmost nonzero entry), clear its column below, move to the next
row and column; the number of pivots of the resulting echelon form
is the rank.

\emph{Inverse computation:} run the algorithm on the block $(A \mid
I_n)$ until the left block becomes $I_n$ (possible iff $A$ is
invertible); the right block is then $A^{-1}$ --- indeed the product
of the elementary matrices used equals $A^{-1}$.
\end{method}

\begin{example}\label{ex:b1:matrices:inverse}
$A = \begin{pmatrix} 1 & 2 \\ 3 & 4 \end{pmatrix}$: reduce $(A \mid
I_2)$:
\[
\begin{pmatrix} 1 & 2 & 1 & 0\\ 3 & 4 & 0 & 1 \end{pmatrix}
\to
\begin{pmatrix} 1 & 2 & 1 & 0\\ 0 & -2 & -3 & 1 \end{pmatrix}
\to
\begin{pmatrix} 1 & 0 & -2 & 1\\ 0 & 1 & \tfrac32 & -\tfrac12
\end{pmatrix},
\]
(operations: $L_2 \leftarrow L_2 - 3L_1$; then $L_1 \leftarrow L_1
+ L_2$, $L_2 \leftarrow -\frac12 L_2$). So $A^{-1} =
\begin{pmatrix} -2 & 1 \\ \tfrac32 & -\tfrac12\end{pmatrix}$.
\emph{Check:} $AA^{-1} = I_2$.
\end{example}

\begin{example}[Rank with a parameter, by rows alone]
\label{ex:b1:matrices:rankparameter}
For $m \in \R$, the rank of
$M_m = \begin{pmatrix} 1 & 1 & m\\ 1 & m & 1\\ m & 1 &
1\end{pmatrix}$. Reduce: $L_2 \leftarrow L_2 - L_1$ and $L_3
\leftarrow L_3 - mL_1$ give the rows
\[
(1,\ 1,\ m), \qquad (0,\ m - 1,\ 1 - m), \qquad
(0,\ 1 - m,\ 1 - m^2).
\]
\emph{Case $m = 1$}: the last two rows vanish --- one pivot,
$\operatorname{rk} M_1 = 1$ (all three original rows were
equal). \emph{Case $m \neq 1$}: scale $L_2$ by $\frac1{m-1}$ and
$L_3$ by $\frac1{1-m}$ to get $(0, 1, -1)$ and $(0, 1, 1 + m)$,
then $L_3 \leftarrow L_3 - L_2 = (0, 0, m + 2)$. If $m = -2$:
two pivots, rank $2$; otherwise three pivots, rank $3$. Summary:
\[
\operatorname{rk} M_m =
\begin{cases}
1 & m = 1,\\
2 & m = -2,\\
3 & \text{otherwise}.
\end{cases}
\]
The same thresholds will drop out of one determinant computation
in \cref{ch:b1:det} (the \omterm{def:b1:poly:def}{polynomial} $-(m+2)(m-1)^2$ of
\cref{exo:b1:det:7}) --- but note what elimination gives that
the determinant does not: the \emph{value} of the rank in the
degenerate cases, not just the fact that it dropped.
\end{example}

\begin{example}[Computing powers]\label{ex:b1:matrices:powers}
$A = \begin{pmatrix} 1 & 1 \\ 0 & 1\end{pmatrix} = I + N$ with $N =
E_{12}$, $N^2 = 0$. Since $I$ and $N$ commute, the binomial theorem
(\cref{prop:b1:structures:binomial}) truncates:
\[
A^k = I + kN = \begin{pmatrix} 1 & k \\ 0 & 1 \end{pmatrix}
\qquad (k \in \N, \text{ and } k \in \Z \text{ using } A^{-1} = I -
N).
\]
\end{example}

\begin{method}[Computing $A^n$: the three routes]
\label{met:b1:matrices:powers}
\begin{enumerate}
  \item \emph{Binomial route}: if $A = \lambda I + N$ with $N$
        nilpotent, the binomial theorem truncates
        (\cref{ex:b1:matrices:powers}, \cref{exo:b1:matrices:5});
        it applies because $\lambda I$ commutes with everything.
  \item \emph{\omterm{def:b1:poly:def}{Polynomial} route}: find a \omterm{def:b1:poly:def}{polynomial} identity
        satisfied by $A$ (in dimension $2$, always $A^2 = sA -
        pI$) and reduce $X^n$ modulo it; the weekend problem
        below builds this route completely.
  \item \emph{Similarity route}: find an invertible $P$ with
        $P^{-1}AP = D$ simple (diagonal, or shift), compute
        $D^n$, and undo: $A^n = P D^n P^{-1}$
        (\cref{thm:b1:matrices:conjugation},
        \cref{ex:b1:matrices:conjugationrun}); the systematic
        search for such $P$ is Year~2's reduction theory.
\end{enumerate}
Whatever the route, check the result on $n = 0, 1, 2$: three
cheap tests that catch nearly every slip.
\end{method}

\begin{remark}[Common pitfalls: the price of non-commutativity]
\label{rem:b1:matrices:pitfalls}
Every identity of scalar algebra whose proof reorders factors
dies in $\mathcal{M}_n(K)$, $n \geq 2$.
\emph{Squares}: $(A + B)^2 = A^2 + AB + BA + B^2$, and the
middle collapses to $2AB$ only if $AB = BA$
(\cref{exo:b1:matrices:1}).
\emph{Powers of products}: $(AB)^k$ is $ABAB\cdots$, not
$A^kB^k$.
\emph{Zero divisors}: $E_{12}E_{12} = 0$ with $E_{12} \neq 0$;
consequently \emph{no cancellation}: $AB = AC$ implies $B = C$
only when $A$ is invertible (multiply by $A^{-1}$ --- on the
correct side).
\emph{\omterm{def:b1:matrices:transpose}{Traces}}: $\operatorname{tr}(AB) = \operatorname{tr}(BA)$
always, but $\operatorname{tr}(AB) \neq
\operatorname{tr}A\operatorname{tr}B$ in general (take $A = B =
I_2$: $2 \neq 4$), and $\operatorname{tr}(ABC) =
\operatorname{tr}(BCA)$ (cyclic) while $\operatorname{tr}(ACB)$
may differ.
\emph{\omterm{def:b1:matrices:transpose}{Transposes} reverse}: $(AB)^{\mathsf T} = B^{\mathsf
T}A^{\mathsf T}$ --- forgetting the reversal is the most common
error in orthogonality computations (\cref{ch:b1:euclid}).
When in doubt, test any claimed identity on $E_{12}$ and
$E_{21}$: the smallest non-commuting pair refutes most false
formulas in one line.
\end{remark}

\begin{remark}[Where the dictionary goes]
\label{rem:b1:matrices:whereused}
The matrix dictionary is used on every remaining page of this
volume: \cref{ch:b1:det} attaches to each square matrix a single
number deciding invertibility, and solves $AX = B$
systematically; \cref{ch:b1:euclid} singles out the matrices
preserving lengths (orthogonal matrices); and in
\cref{ch:b1:multivar}, the second-order behavior of a function of
two variables is a symmetric $2 \times 2$ matrix. The \omterm{def:b1:matrices:transpose}{trace},
introduced above almost in passing, becomes a powerful invariant:
\cref{exo:b1:matrices:6,exo:b1:matrices:8} give a first taste,
and the Year~2 volume builds eigenvalue theory on it. The weekend
problem develops the other workhorse: \omterm{def:b1:poly:def}{polynomial} identities
satisfied by a matrix, which turn the computation of $A^n$ into
a two-term \omterm{pb:b1:matrices:1}{linear recurrence}.
\end{remark}

\begin{remark}[Perspectives inside Book 3]
\label{rem:b1:matrices:perspectives}
Three matrix families introduced here have appointments later in
this volume. \emph{Symmetric matrices}
(\cref{ex:b1:matrices:symsplit}) carry the second-order data of
functions of two variables: the Monge test of
\cref{ch:b1:multivar} is a \omterm{def:b1:logic:statement}{statement} about the sign behavior of
a symmetric $2\times2$ matrix, and its determinant $rt - s^2$
is computed by \cref{ch:b1:det}'s machinery. \emph{Orthogonal
matrices} ($A^{\mathsf T}A = I$) are the isometries of
\cref{ch:b1:euclid}, where the \omterm{def:b1:matrices:transpose}{transpose} finally acquires its
geometric meaning: it is the algebraic shadow of the inner
product. \emph{Invertible matrices} meet their practical test
in \cref{ch:b1:det} --- one number, $\det A \neq 0$ --- closing
the search this chapter began with row reduction. \omterm{def:b1:matrices:transpose}{Trace} and
determinant then travel as the invariant pair $(s, p)$ of the
weekend problem, all the way to the eigenvalue theory of
Year~2.
\end{remark}

\section{Exercises}

\begin{exercise}[$\star$]\label{exo:b1:matrices:1}
Let $A = \begin{pmatrix} 1 & 2 \\ 0 & 1 \end{pmatrix}$ and $B =
\begin{pmatrix} 0 & 1 \\ 1 & 0\end{pmatrix}$. Compute $AB$, $BA$,
$A^2 - B^2$ and $(A+B)(A-B)$; explain why the last two differ.
\end{exercise}

\begin{exercise}[$\star$]\label{exo:b1:matrices:2}
Compute the rank of
\[
M = \begin{pmatrix}
1 & 2 & 3\\
2 & 4 & 6\\
1 & 1 & 1
\end{pmatrix},
\qquad
N = \begin{pmatrix}
1 & 1 & 0 & 2\\
0 & 1 & 1 & 1\\
1 & 2 & 1 & 3
\end{pmatrix}.
\]
\end{exercise}

\begin{exercise}[$\star$]\label{exo:b1:matrices:3}
Invert, by row reduction, $A = \begin{pmatrix} 1 & 0 & 1\\ 2 & 1 &
1\\ 1 & 1 & 1 \end{pmatrix}$, and check on one product.
\end{exercise}

\begin{exercise}[$\star$]\label{exo:b1:matrices:4}
Write the matrix, in the canonical \omterm{def:b1:vspaces:free}{basis} of $\R_2[X]$, of the
endomorphism $u(P) = P(X + 1)$. Explain, without computation, why
it is invertible, and give the matrix of $u^{-1}$.
\end{exercise}

\begin{exercise}[$\star\star$]\label{exo:b1:matrices:5}
Let $A = \begin{pmatrix} 2 & 1 \\ 0 & 2\end{pmatrix}$. Write $A =
2I + N$, compute $N^2$, and deduce $A^k$ for all $k \in \N$ by the
binomial theorem.
\end{exercise}

\begin{exercise}[$\star\star$]\label{exo:b1:matrices:6}
Prove that there are no matrices $A, B \in \mathcal{M}_n(K)$ (with
$K = \R$ or $\C$) such that $AB - BA = I_n$. \emph{(Take \omterm{def:b1:matrices:transpose}{traces}.)}
\end{exercise}

\begin{exercise}[$\star\star$]\label{exo:b1:matrices:7}
A matrix $A$ is \emph{nilpotent} when $A^m = 0$ for some $m$. Prove
that $I - A$ is then invertible, with
\[
(I - A)^{-1} = I + A + A^2 + \dots + A^{m-1} .
\]
Application: invert $\begin{pmatrix} 1 & 2 & 3\\ 0 & 1 & 2\\ 0 & 0 &
1\end{pmatrix}$.
\end{exercise}

\begin{exercise}[$\star\star$]\label{exo:b1:matrices:8}
Let $A \in \mathcal{M}_n(\R)$ satisfy $A^2 = A$ (idempotent). Prove
that $\operatorname{tr} A = \operatorname{rk} A$. \emph{(Interpret
$A$ as a \omterm{def:b1:linmaps:projection}{projection} and choose an adapted \omterm{def:b1:vspaces:free}{basis};
\cref{thm:b1:matrices:conjugation} says the \omterm{def:b1:matrices:transpose}{trace} is
basis-independent since $\operatorname{tr}(P^{-1}MP) =
\operatorname{tr} M$.)}
\end{exercise}

\begin{exercise}[$\star\star\star$]\label{exo:b1:matrices:9}
Let $J \in \mathcal{M}_n(\R)$ be the all-ones matrix. Compute $J^2$,
and deduce, for $a, b \in \R$, the condition of invertibility of $M
= aI + bJ$ together with $M^{-1}$ \emph{(look for an inverse of the
same form $\alpha I + \beta J$)}.
\end{exercise}

\begin{exercise}[$\star\star\star$]\label{exo:b1:matrices:10}
(Rank inequalities) For $A, B \in \mathcal{M}_n(K)$, prove
\[
\operatorname{rk}(A + B) \leq \operatorname{rk} A +
\operatorname{rk} B,
\qquad
\operatorname{rk}(AB) \geq \operatorname{rk} A + \operatorname{rk}
B - n .
\]
\emph{(For the second --- Sylvester's inequality --- apply
rank--nullity to the restriction of the \omterm{def:b1:logic:map}{map} of $A$ to
$\operatorname{im} B$.)}
\end{exercise}

\begin{exercise}[$\star\star$]\label{exo:b1:matrices:11}
Let $D = \operatorname{diag}(d_1, \dots, d_n)$ with the $d_i$
\emph{pairwise distinct}.
\begin{enumerate}
  \item Prove that a matrix $A$ commutes with $D$ if and only if
        $A$ is diagonal. \emph{(Compare the $(i,j)$ entries of
        $AD$ and $DA$.)}
  \item Deduce the \emph{center} of $\mathcal{M}_n(K)$: the
        matrices commuting with \emph{every} matrix are exactly
        the scalar matrices $\lambda I_n$. \emph{(Test against
        $D$, then against the matrices $E_{ij}$.)}
\end{enumerate}
\end{exercise}

\begin{exercise}[$\star\star\star$]\label{exo:b1:matrices:12}
(Rank-one matrices) Let $A \in \mathcal{M}_n(K)$, $A \neq 0$.
\begin{enumerate}
  \item Prove that $\operatorname{rk} A = 1$ if and only if $A =
        CL$ for a nonzero column $C \in \mathcal{M}_{n,1}$ and a
        nonzero row $L \in \mathcal{M}_{1,n}$.
  \item For such an $A$, prove $A^2 = (\operatorname{tr} A)\,A$;
        deduce that a rank-one matrix is nilpotent if and only if
        its \omterm{def:b1:matrices:transpose}{trace} is zero.
  \item If $\operatorname{tr} A \neq -1$, prove that $I_n + A$ is
        invertible with
        \[
        (I_n + A)^{-1} = I_n - \frac{1}{1 + \operatorname{tr}
        A}\,A ,
        \]
        and that $I_n + A$ is \emph{not} invertible when
        $\operatorname{tr} A = -1$. \emph{(Find a vector killed
        by $I_n + A$.)}
\end{enumerate}
\end{exercise}

\section{Problem: powers of a matrix by polynomial division}

\begin{problem}\label{pb:b1:matrices:1}
Computing $A^{100}$ entry by entry is hopeless; computing it
through a \omterm{def:b1:poly:def}{polynomial} identity satisfied by $A$ is three lines.
This problem builds the method from scratch: euclidean division
of $X^n$, the identity $A^2 - sA + pI = 0$ verified by every $2
\times 2$ matrix (the Cayley--Hamilton theorem in dimension
$2$), and the dictionary between matrix powers and \omterm{pb:b1:matrices:1}{linear
recurrences} --- with the \omterm{pb:b1:matrices:1}{Fibonacci numbers} as running example.
\index{Cayley--Hamilton theorem (dimension 2)}
\index{Fibonacci numbers}\index{linear recurrence}

\medskip
\noindent\textbf{Part I --- The remainder calculus.} Fix $s, p
\in K$ and $D = X^2 - sX + p$.
\begin{enumerate}
  \item Justify that for each $n \in \N$ there are unique $Q_n
        \in K[X]$ and $(a_n, b_n) \in K^2$ with
        \[
        X^n = Q_n\,D + a_n X + b_n ,
        \]
        and compute $(a_0, b_0)$ and $(a_1, b_1)$.
  \item Multiplying by $X$ and dividing again, establish the
        recurrences
        \[
        a_{n+1} = s\,a_n + b_n,
        \qquad
        b_{n+1} = -p\,a_n ,
        \]
        and deduce $a_{n+2} = s\,a_{n+1} - p\,a_n$: the
        coefficient sequence obeys the \omterm{pb:b1:matrices:1}{linear recurrence} attached
        to $D$.
  \item Suppose $D$ has two distinct roots $\lambda \neq \mu$.
        Evaluating the division identity, prove
        \[
        a_n = \frac{\lambda^n - \mu^n}{\lambda - \mu},
        \qquad
        b_n = \frac{\lambda\mu^n - \mu\lambda^n}{\lambda - \mu} .
        \]
  \item Suppose $D = (X - \lambda)^2$. Using the \omterm{def:b1:derivative:def}{derivative} of
        the division identity, prove $a_n = n\lambda^{n-1}$ and
        $b_n = (1 - n)\lambda^{n}$.
  \item Show that substituting a fixed matrix $M \in
        \mathcal{M}_k(K)$ into \omterm{def:b1:poly:def}{polynomials} respects sums and
        products: $(PQ)(M) = P(M)\,Q(M)$. Deduce that if $D(M) =
        0$, then
        \[
        M^n = a_n\,M + b_n\,I \qquad (n \in \N).
        \]
\end{enumerate}

\medskip
\noindent\textbf{Part II --- Dimension 2: \omterm{def:b1:matrices:transpose}{trace}, determinant
number, Cayley--Hamilton.} For $A = \begin{pmatrix} a & b\\ c &
d\end{pmatrix}$ set $s = a + d = \operatorname{tr} A$ and $p = ad
- bc$ (the number that \cref{ch:b1:det} will name the
determinant).
\begin{enumerate}[resume]
  \item Verify by direct computation the \emph{Cayley--Hamilton
        identity in dimension $2$}:
        \[
        A^2 - s\,A + p\,I_2 = 0 .
        \]
  \item Prove by direct expansion that $p$ is multiplicative:
        with obvious notation, $p(AB) = p(A)\,p(B)$. Then show:
        $A$ is invertible if and only if $p \neq 0$, in which
        case
        \[
        A^{-1} = \frac1p\,\bigl(s\,I_2 - A\bigr).
        \]
  \item Let $A = \begin{pmatrix} 1 & 1\\ 0 & 2\end{pmatrix}$.
        Compute $s$, $p$, the roots of $D$, and deduce a \omterm{def:b1:topology:closed}{closed}
        formula for $A^n$; check it against a direct computation
        of $A^2$.
  \item Let $A = \begin{pmatrix} 3 & 1\\ -1 & 1\end{pmatrix}$.
        Show that $D$ has a double root and compute $A^n$; check
        at $n = 2$.
  \item Let $F = \begin{pmatrix} 1 & 1\\ 1 & 0\end{pmatrix}$ and
        define the \omterm{pb:b1:matrices:1}{Fibonacci numbers} by $F_0 = 0$, $F_1 = 1$,
        $F_{n+2} = F_{n+1} + F_n$. Prove
        \[
        F^n = \begin{pmatrix} F_{n+1} & F_n\\ F_n &
        F_{n-1}\end{pmatrix} \quad (n \geq 1),
        \]
        deduce Binet's formula $F_n = \dfrac{\varphi^n -
        \psi^n}{\sqrt5}$ where $\varphi = \frac{1 + \sqrt5}2$,
        $\psi = \frac{1 - \sqrt5}2$, and, using question 7,
        Cassini's identity $F_{n+1}F_{n-1} - F_n^2 = (-1)^n$.
\end{enumerate}

\medskip
\noindent\textbf{Part III --- \omterm{pb:b1:matrices:1}{Linear recurrences}, structurally.}
Fix $s, p \in K$ with $p \neq 0$, and let $E_D$ be the \omterm{def:b1:logic:sets}{set} of
sequences with $u_{n+2} = s\,u_{n+1} - p\,u_n$ for all $n$.
\begin{enumerate}[resume]
  \item Show that $E_D$ is a \omterm{def:b1:vspaces:def}{vector space} of dimension $2$
        (adapt \cref{exo:b1:findim:10}).
  \item Show that the sequence $(a_n)$ of Part I is the element
        of $E_D$ with initial values $0, 1$, and that every $u
        \in E_D$ satisfies
        \[
        u_n = u_1\,a_n + u_0\,b_n \qquad (n \in \N),
        \]
        with $(b_n)$ as in Part I: the division remainders solve
        \emph{all} recurrences at once.
  \item If $\lambda \neq \mu$ are the roots of $D$, show that
        $\bigl((\lambda^n), (\mu^n)\bigr)$ is a \omterm{def:b1:vspaces:free}{basis} of $E_D$;
        if $D = (X-\lambda)^2$ with $\lambda \neq 0$, show that
        $\bigl((\lambda^n), (n\lambda^n)\bigr)$ is one.
  \item Solve completely: $u_{n+2} = u_{n+1} + 6u_n$, $u_0 = 1$,
        $u_1 = 8$; check the answer on $u_2$ and $u_3$.
  \item Let $C = \begin{pmatrix} 0 & 1\\ -p & s\end{pmatrix}$
        (the \emph{companion matrix} of $D$). Show that
        \[
        \begin{pmatrix} u_{n}\\ u_{n+1}\end{pmatrix}
        = C^n \begin{pmatrix} u_0\\ u_1\end{pmatrix}
        \quad (u \in E_D),
        \]
        and that $\operatorname{tr} C = s$ and $p(C) = p$: the
        recurrence and the matrix carry the same \omterm{def:b1:poly:def}{polynomial} $D$.
\end{enumerate}

\medskip
\noindent\textbf{Part IV --- Degree three.} Let $D_3 = X^3 -
\alpha X^2 - \beta X - \gamma$ and
\[
C_3 = \begin{pmatrix} 0 & 1 & 0\\ 0 & 0 & 1\\ \gamma & \beta &
\alpha \end{pmatrix}.
\]
\begin{enumerate}[resume]
  \item Show that $D_3(C_3) = 0$. \emph{(Compute the \omterm{def:b1:linmaps:kerim}{images} of
        the canonical \omterm{def:b1:vspaces:free}{basis} vectors under powers of $C_3$: the
        \omterm{def:b1:logic:map}{map} of $C_3$ sends $e_1 \mapsto \dots \mapsto$ a
        combination forced by the last row.)}
  \item Show that if $D_3$ has three distinct roots $\lambda_1,
        \lambda_2, \lambda_3$, the remainder $R_n$ of $X^n$
        divided by $D_3$ is the \emph{Lagrange interpolant} of
        the values $\lambda_i^n$ at the nodes $\lambda_i$
        (\cref{thm:b1:poly:lagrange}); deduce that every entry
        of $C_3^{\,n}$ is a fixed linear combination of
        $\lambda_1^n, \lambda_2^n, \lambda_3^n$.
  \item Solve: $u_{n+3} = 2u_{n+2} + u_{n+1} - 2u_n$ with $u_0 =
        0$, $u_1 = 1$, $u_2 = 1$. \emph{(Factor $D_3 = (X - 1)(X
        + 1)(X - 2)$.)} Check on $u_3$.
  \item Compute the remainder of $X^n$ modulo $(X -
        \lambda)^3$ \emph{(Taylor expansion of $X^n$ at
        $\lambda$)}, and deduce a formula for $(\lambda I +
        N)^n$ when $N^3 = 0$ and $N$ commutes with everything in
        sight; check it against the binomial theorem.
  \item Show that for $D_3$ with distinct roots, the general
        solution of the order-$3$ recurrence is $u_n = c_1
        \lambda_1^n + c_2\lambda_2^n + c_3\lambda_3^n$: prove
        that the three geometric sequences form a \omterm{def:b1:vspaces:free}{basis} of the
        solution space. \emph{(For freeness, evaluate a null
        combination at $n = 0, 1, 2$ and recognize an
        interpolation system at the distinct nodes
        $\lambda_i$.)}
\end{enumerate}

\medskip
\noindent\textbf{Part V --- Fibonacci dividends, and synthesis.}
\begin{enumerate}[resume]
  \item Prove $F_1 + F_2 + \dots + F_n = F_{n+2} - 1$.
  \item From $F^{m+n} = F^m F^n$, derive the addition formula
        \[
        F_{m+n} = F_{m+1}F_n + F_m F_{n-1},
        \]
        and deduce $F_{2n} = F_n(F_{n+1} + F_{n-1})$.
  \item Prove that $F_n$ is the nearest integer to
        $\varphi^n/\sqrt5$ for every $n \geq 0$.
  \item Let $t_n = \operatorname{tr}(F^n) = F_{n+1} + F_{n-1}$
        (the \emph{Lucas numbers} $L_n$). Show $t_{n+2} =
        t_{n+1} + t_n$, $t_1 = 1$, $t_2 = 3$, that $L_n =
        \varphi^n + \psi^n$, and recover $F_{2n} = F_n L_n$.
  \item Synthesis, in four sentences: why the powers of a $2
        \times 2$ matrix live in the plane
        $\operatorname{Vect}(I, A)$ of $\mathcal{M}_2(K)$ (which
        dimension argument guarantees a quadratic identity, and
        which explicit identity Part II produced); how euclidean
        division converts exponentiation into a two-term
        recurrence; which \omterm{def:b1:logic:statement}{statement} of this problem is the $n =
        2$ case of a theorem valid in all dimensions (name it,
        and say where it is proved in this \omterm{def:b1:series:def}{series}); and what the
        companion matrix construction adds to the picture.
\end{enumerate}
\end{problem}
